\section{怎样掌握和运用成语}

成语是作为完整的意义单位来运用的，定型的一种固定词组，它在句子中的作用往往相当于一个词。它有两个特点：一是结构上定型化；二是意义上整体化。汉语的成语非常精练、形象而又精辟，具有强大的生命力和深刻的表现力，因此在社会生活中得到广泛的应用。我们要培养自己具有正确地、自如地运用汉语词语的能力，也必得掌握好成语才行。问题是，汉语成语成千上万，怎样才能掌握并能正确地运用它呢？

1. 注意掌握成语的结构关系。

汉语成语的普遍格式是四字格，四字格化是成语产生和发展的总趋势。这种四字格成语，节奏鲜明，声律铿锵，音调和谐，形式均衡，富有音乐美和旋律感，因此读来琅琅上口，容易记忆。

这种四字格成语的结构关系与一般词组的结构关系大致相同，有主谓关系、联合关系、偏正关系、动宾关系、述补关系，等等。例如：

主谓关系：众星捧月、塞翁失马、鹤立鸡群、言过其实、毛遂自荐、敝帚自珍（一主一语）；声色俱厉，鱼龙混杂，鸡犬不宁，青黄不接，
休戚相关，经纬万端（二主一谓）；似是而非，不寒而栗，听而不闻，临危不惧，好为人师，恬不知耻（只有谩语）。

联合关系：海枯石烂，家喻户晓，貌合神离，星罗棋布，唇亡齿寒，字斟句酌（两个主谓式联合）；良师益友，雄才大略，大惊小怪，赤手空拳，
千锤百炼，远走高飞（两个偏正式联合）；扶危济困，改弦易辙，藏垢纳污，背井离
乡，开宗明义，驾轻就熟（两个动宾式联
合）；

光明磊落，简明扼要，慷慨激昂，错综复
杂，汹涌澎湃，富丽堂皇（两个双音词联
合）。

偏正关系：一丘之貉，锦绣前程，过眼云烟，乌合之众，巍然屹立，惴惴不安，谆谆告诫，豁然开朗，无耻之尤，幡然悔悟。

动宾关系：叱咤风云，崭露头角，墨守成规，震撼人心，重整旗鼓，别开生面，大放厥词。

述补关系：毁于一旦，掉以轻心，瞠乎其后，置之度外，嗤之以鼻，危在旦夕，俗不可耐。

兼语式：点铁成金、引狼入室、发人深省、调虎离山、震耳欲聋、挥汗成雨、引人入胜。

连谓式：哗众取宠，含沙射影，见微知著，打草惊蛇，守株待兔，削足适履，见利忘义，缘木求鱼。

\newpage
\noindent
解甲归田、挂一漏万

我们在学习成语时，从了解一个成语的结构关系入手，可以利用同义词和反义词在成语中的运用，来帮助我们记住成语和理解成语中不懂的词义。这一点，在学习那些结构是联合关系的成语时是很有用的。

2. 注意掌握成语的确切的含义。

成语虽然是由词组组成的，但是只有一部分成语可以从词面上加以解释。因此，理解成语的意义，一方面要弄清字义，更重要的是要领会它的整体意义——有的是引申意义，有的是比喻意义，而决不能望文生义。特别对一些由古书中的故事或词句概括而成的成语，一定要弄明白它们的来源及其特定的意义，才能正确地理解它们的整体意义。

例如：“废寝忘食”、“顾此失彼”、“牢不可破”、“顾名思义”、“外强中干”、“名存实亡”这一类成语，根据字义和彼此之间的结构关系，可以大体了解它们整体的意思。而象“黄粱一梦”、“满城风雨”、“胸有成竹”、“四面楚歌”、“水落石出”、“青出于蓝”、“为虎作伥”、“破釜沉舟”等等成语，它的整体意义和构成这一成语的字义，是不能等同的，而要从它们的比喻义来理解整体的含义。这一类来源于古诗文的成语，决不能局限于它们的字面含义去进行解释，也就是说，不能从各个字的意义的累加与总和中去了解这一类成语的整体意义。例如“黄粱一梦”这个成语，是对一个故事的概括。据唐人沈既济《枕中记》载：卢生在邯郸客店中昼寝入梦，历尽富贵荣华；梦醒，主人烧黄粱饭还未熟。因此对这个成语的正确理解是：在煮一锅黄米饭的短暂时间里做了一场好梦（字面的意义）；
\newpage
\noindent
比喻虚幻梦想，欲望一场空（比喻的意义），而说出它的比喻义是解释这个成语的根本要求。“黄粱一梦”，又作“黄粱美梦”、“一枕黄粱”，都是同出一源的同义成语。再如“胸有成竹”这个成语，是对一句古文的概括。宋人苏轼《文与可画筼筜谷偃竹记》有这样一句话：“故画竹，必先得成竹于胸中。”因此，对这个成语的正确理解是：在画竹之前，一定要在心里构思好现成的完整的竹子（字面的意义）；比喻处理事情先有主意，先有成熟的考虑（比喻的意义），而解释出它的比喻义，是对解释这个成语的根本要求。“胸有成竹”，又作“成竹在胸”，都是同出一源的同义成语。如果只按字面解释而望文生义，是不行的，甚至会闹出笑话来。

我们在解释成语时，要特别注意辨析其中关键的字，有时还要借助所学的古代汉语有关知识来理解，因为有不少成语是从古代文献中的词句或故事概括而成的。例如：

狐假虎威：假，借着。比喻倚仗别人的势力来欺压人。

首鼠两端：首鼠，犹豫不决，欲进又退的样子，形容瞻前顾后、迟疑不决的样子。

弃甲曳兵：兵，兵器。丢掉铠甲，拖着兵器，形容打败仗逃跑时的狼狈相。

衣锦还乡：衣（yì），穿。旧时指富贵以后回到家乡，向亲友夸耀。

釜底抽薪：釜，锅。从锅底下抽掉柴火，比喻从根本上解决问题。

3. 注意分清成语的感情色彩。

成语同一-般词语一样，有褒义、贬义和中性之分。我们

\newpage
\noindent
在学习和运用成语时，一定要根据句子所表达的感情色彩，去辨析和选择恰当的成语。例如：

带褒义的：扬眉吐气、无微不至、深思熟虑、画龙点睛、自食其力、再接再厉、红光满面、苦心孤诣、有口皆碑、柳暗花明

带贬义的：趾高气扬、无所不至、处心积虑、画蛇添足、自食其果、变本加厉、面红耳赤、挖空心思、螳臂当车、日暮途穷

如果不分清成语的褒贬色彩，在表示赞扬或肯定意思的场合误用了贬义成语，或者在表示贬斥或否定意思的场合误用了褒义成语，都是不行的。例如：“打倒了四人帮”，人民群众趾高气扬了。”这个句子不应该用贬义成语“趾高气扬”，而应该用褒义成语“扬眉吐气”。又如：“在‘文革’期间，那些打砸抢分子胡作非为，无微不至。”这个句子不应该用褒义成语“无微不至”，而应该用贬义成语“无所不至”。究竟是该用褒义成语，还是该用贬义成语，可以从前后词语的配合与整个句子的感情色彩辨析出来。据此，我们就能避免或者改正由于所用的成语不合句意所表达的褒贬色彩而造成的病句。例如：

① 在演讲比赛中，王莉信口开河，滔滔不绝，终于获得了一等奖。
（从“获得了一等奖”来看，应将贬义成语“信口开河”改为褒义成语“口若悬河”。）

② 这个坏家伙由于老谋深算，善于伪装，一直隐藏到现在。
（既然说的是“坏家伙”，就不应用褒义成语“老谋深

\newpage
\noindent
算”，而应改用贬义成语“老奸巨猾”。）

③ 他同我们认识不久，就肆无忌惮地和我们谈起他的经历和对生活的看法，看来他是一个心地坦率的人。
（既然“看来……心地坦率”，就不应用贬义成语“肆无忌惮”，而应改用中性成语“百无禁忌”。）

④ 同学们上课时聚精会神，专心听讲；在操场上生龙活虎，杀气腾腾。
（全句都带褒意色彩，所以应将贬意成语“杀气腾腾”改为褒意成语“热气腾腾”。）

⑤ 在和反动派斗争时，必须善于利用矛盾，两面三刀，因为敌人从来就不是铁板一块的。
（因为是从人民的角度说的，所以必须将贬义成语“两面三刀”改用褒义成语“各个击破”。）

4. 注意掌握成语的正确写法和读音。

写错或读错成语，大多是由于没有真正理解成语的含义造成的。特别是对于成语中的文言词语必须按照古代汉语的词义去理解，而不能用现代汉语的词义去理解，如果理解了，就可能写错字，读错音，使整个成语不知所云了。例如：

数见不鲜（数：shuò，屡次；不能读成 shǔ。鲜：xiān，新奇；不能读为 xiǎn。）

汗流浃背（浃：jiā，湿透；不能写成“夹”，也不能读为 jiǎ。）

居心叵测（叵：pǒ，不可；不能写成“臣”，也不能读成 jù。）

日薄西山（薄：bó，靠近；不能写成“薄”，也不能读

\newpage
\noindent
成 bù。）

生死不渝（渝：yú，改变；不能写成“逾”，也不能读成 yù。）

怙恶不悛（怙：hù，仗恃，坚持；不能写成“估”，也不能读成 gǔ。悛：quān，改悔；不能写作“俊”，也不能读作 jùn。）

鳞次栉比（栉：zhì，梳篦的总称；不能写成“节”，也不能读作 jiē。）

稳操左券（券：quàn，古代契约中左右两联的一联，作为索偿的凭证；不能写成“卷”，也不能读作 juàn。）

负隅顽抗（隅：yú，同“嵎”，山势弯曲险阻的地方；不能写成“偶”，也不能读作 ǒu。）

一暴十寒（暴：pù，晒；不能写成“爆”，也不能读作 bào。）

从上面这些错写错误的例子里，我们可以看出，要正确掌握和运用一个成语，必须弄懂这个成语中的一两个关键字，否则就会读错或写错。比如“严惩不贷”的“贷”字，经常被误写成“代”、“待”，就是因为没弄清“贷”是“饶恕、宽容”的意思；而“责无旁贷”的“贷”字，也经常被误写成“代”、“待”，则是因为没弄清这个“贷”字是“推卸”的意思。再如“变本加厉”的“厉”字，经常被误写成“利”、“励”甚至“历”，就是因为没弄清整个成语的含义是形容情况比原来更加严重，“厉”是“厉害、猛烈”的意思。又如“不可思议”的“议”字，经常被误写成“意”或“义”，也是由于不了解整个成语的含义是“不可

\newpage
\noindent
想象，不可理解”的意思（这个成语原来是佛教用语，含有神秘奥妙的意思），于是就望文生义地错写成“不可思意”或“不可思义”了。此外如“道貌岸（不能写成“暗”）然”、“十恶不赦（不能写成“舍”、“赧”）”、“草菅（不能写成“管”）人命”、“如火如荼（不能写成“茶”）”、“戮（不能写成“戳”）力同心”、“鬼鬼祟祟（不能写成“崇崇”）”、“恬（不能写成“括”）不知耻”、“怨天尤（不能写成“由”）人”、“为虎作伥（不能写成“胀”）”、“无事（不能写成“是”）生非”，等等。只要我们留心辨析这类容易讲错、写错或读错的成语，并且持之以恒，就可以少犯或者不犯这类的错误了。

5. 注意成语的规范性。

作为固定词组的成语，有它紧密的结构形式和完整的概念意义，因此不能任意改换成语的字眼和变更成语的结构形式。例如：
“遍体鳞伤”不能改写成“全体鳞伤”；
“挂一漏万”不能改写成“挂一漏亿”；
“好为人师”不能改写成“好当人师”……

这些成语，改换了意义相近的字，整个成语的意思也大体相同，但是它破坏了成语的极强的稳定性，违反了人们约定俗成的习惯用法。至于要改换成语中字音或字形相近的字，更是不能允许的，因为那样就使成语的意思改变了或者不能理解了。例如：“置之度外”不能写成“置之肚外”，“一筹莫展”不能写成“一愁莫展”，等等；“心力交瘁”不能写成“心力交瘁”，“据为己有”不能写成“据为己有”，“休戚相关”不能写成“休戚相关”，“漠不关心”不能写

\newpage
\noindent
成“莫不关心”，等等。

成语的结构形式也不能随便改变，因为随便改变，不符合人们习惯的用法，有的则改变了原来成语的意思，有的甚至不可理解了。例如：“有声有色”写成“有色有声”、“耳提面命”写成“面命耳提”、“诚惶诚恐”写成“诚恐诚惶”、“狐群狗党”写成“狗党狐群”、“错综复杂”写成“复杂错综”，“十全十美”写成“十美十全”，等等，大体意思虽然没变，但是人们约定俗成，不习惯这样用。再如：“前仆后继”写成“后继前仆”、“老马识途”写成“识途老马”、“痴人说梦”写成“说梦痴人”，这样一改，就不能反映原来成语的那种意思了。又如“过河拆桥”写成“过桥拆河”，“驾轻就熟”（驾轻车走熟路）写成“驾熟就轻”，若这样一改，那简直不知所云了。

当然，有少数成语在发展过程中，变动了个别字眼或结构形式，但是这种变动已为社会所公认，符合人们习用的约定俗成的原则，就不在上述范围之内了。如“前仆后继”，现在也可以写成“前赴后继”，完全赋予了新的含义，并为人们所习用，当然不在此例。此外，由于时代不同，人们的认识不同，有些成语的含义也有了显著的变化。如“一团和气”本来是指人的态度“和蔼可亲”，带有褒义，现在把无原则的团结叫做“一团和气”，转化为贬义了。又如“矫枉过正”本来指为要把弯曲的东西扭直，结果又歪向了另一方，比喻纠正错误而超过了应有的限度，现在则指必须用群众的革命方法，而不能用改良方法去结束剥削阶级统治的旧秩序，转化为褒义了。我们作为当代人，在理解和运用成语时，一般应该按照现在大家习用的褒贬色彩去讲、去用才好。